\documentclass[11pt]{article}

\usepackage[T1]{fontenc}
\usepackage[margin=1in]{geometry}
\usepackage{booktabs}
\usepackage{tabularx}
\usepackage{listings}
\usepackage{xcolor}

\lstset{
  basicstyle=\ttfamily\small,
  breaklines=true,
  frame=single,
  columns=fullflexible
}

\title{ADR-0001: Integration of Enterprise Network Capture and Alert Planning into the LexLean Wireshark Rewrite}
\author{Wireshark LexLean Rewrite Working Group}
\date{2026-11-23}

\begin{document}
\maketitle

\begin{abstract}
This decision record governs the integration of
\texttt{Network Capture and Alert Planning.txt} into the LexLean 1.2 Wireshark rewrite located at
\texttt{/run/media/multiplic/AGI/prismagi/mount/SUBSYSTEMS/production}. The source document is a
rich operational architecture narrative covering packet acquisition, passive sensors, Kafka
buffering, schema normalization, file-server threat analytics, and SOAR response. This ADR decides
to integrate its modelable claims as new closed-lexicon, kernel-verified LexLean modules rather than
as copied prose, operational configuration, or a claim of equivalence with deployed infrastructure.
\end{abstract}

\section{Context}

The existing project is a bounded formal rewrite in the style of \texttt{kuramoto-prism} and
\texttt{o3de-prism}. Before this change it contained nine verified modules covering pcap/pcapng
capture limits, block alignment, dissector fields, display filters, cross-module invariants, PRISM
architecture/quality/security facets, and top-level integration.

The planning document introduces five additional concern areas:

\begin{enumerate}
  \item a four-tier high-throughput capture pipeline: traffic acquisition, passive sensor
        processing, decoupled transport buffering, and ingestion/storage;
  \item schema normalization across native Zeek fields, ECS, OCSF, and Google UDM, including the
        OCSF identity $\mathrm{type\_uid} = \mathrm{class\_uid}\times 100 + \mathrm{activity\_id}$;
  \item a four-phase file-server threat lifecycle and four detection families: remote payload
        creation on administrative shares, automated share enumeration, rapid ransomware
        rename/overwrite, and anomalous mass read/exfiltration;
  \item baseline filtering, contextual scoring, alert severity, and network/endpoint correlation;
  \item SOAR response playbooks mapped to the resulting alert class.
\end{enumerate}

The implementation plan preceding this ADR is
\texttt{docs/PLAN-network-capture-alert-integration.md}.

\section{Decision Drivers}

\begin{enumerate}
  \item \textbf{Bounded formal fidelity.} Model the assertions and thresholds that can be stated
        finitely and proved, without pretending to reproduce external systems.
  \item \textbf{Closed-lexicon discipline.} Every new module must use the project's closed title
        lexicon and follow LexLean 1.2 canonical formatting.
  \item \textbf{Kernel verification.} Every theorem must be discharged by the Lean 4 kernel with no
        sorries and no dynamic tactic axioms.
  \item \textbf{Requirements traceability.} Each major planning area must map to named LexLean
        declarations and to a verification command.
  \item \textbf{Additive integration.} Existing Wireshark modules and the original
        \texttt{../wireshark} source tree remain unchanged.
  \item \textbf{Release reproducibility.} The complete expanded project must pass \texttt{prismpm
        check} and publish a new \texttt{prismpm build}.
\end{enumerate}

\section{Decision}

We add five new LexLean modules under \texttt{src/Enterprise/}, extend the top-level
\texttt{Main} module, add five title entries to \texttt{language/wireshark.titles}, and extend the
sorted \texttt{entrypoints} array in \texttt{lexlean.toml}. The source planning document remains an
input artifact; it is not copied into the formal source and no external deployment is performed.

\subsection{Module Set}

\begin{table}[htbp]
\centering
\small
\begin{tabularx}{\textwidth}{l X}
\toprule
\textbf{Module} & \textbf{Modeled responsibility} \\
\midrule
\texttt{Enterprise.CapturePipeline}
  & Mirroring encodings (physical TAP, SPAN, VXLAN/UDP, GENEVE/UDP), four-tier ordering, tunnel
    termination, \texttt{SO\_REUSEPORT}-style worker balancing, and buffered transport. \\
\texttt{Enterprise.Normalization}
  & Schema identifiers, canonical source/destination/file field roles, OCSF class and activity
    codes, and the OCSF \texttt{type\_uid} formula. \\
\texttt{Enterprise.FileServerThreats}
  & Attack phases, share kinds, file actions, extension classes, and the four rule predicates with
    their documented thresholds. \\
\texttt{Enterprise.AlertOperations}
  & Baseline suppression, contextual weights, score thresholding, alert severity, endpoint
    correlation, and response action codes. \\
\texttt{Enterprise.Invariants}
  & Cross-module consistency theorems linking pipeline, normalization, detection, and alerting. \\
\bottomrule
\end{tabularx}
\caption{New formal modules added by this ADR.}
\end{table}

\subsection{Representative Formalized Constants}

\begin{table}[htbp]
\centering
\small
\begin{tabular}{lll}
\toprule
\textbf{Concept} & \textbf{Value} & \textbf{Source intent} \\
\midrule
OCSF File System Activity class UID & 1001 & Class 1001 \\
OCSF Network Activity class UID & 4001 & Class 4001 \\
File creation activity ID & 1 & $\texttt{type\_uid}=100101$ \\
File renaming activity ID & 5 & $\texttt{type\_uid}=100105$ \\
Share enumeration threshold & $\geq 15$ distinct targets & Rule 2, ten-second window \\
Ransomware event threshold & $\geq 50$ events & Rule 3, one-minute window \\
Exfiltration read threshold & $\geq 500$ events & Rule 4, five-minute window \\
\bottomrule
\end{tabular}
\caption{Constants preserved from the planning document.}
\end{table}

\section{Requirements Traceability}

\begin{table}[htbp]
\centering
\small
\begin{tabularx}{\textwidth}{p{3.6cm} X l}
\toprule
\textbf{Planning area} & \textbf{Formal claim} & \textbf{Declaration family} \\
\midrule
Four-tier pipeline & Tier indices are strictly ordered and buffering is present in the default
  pipeline. & \texttt{tierIndex}, \texttt{isPipelineWellFormed},
  \texttt{defaultEnterprisePipeline} \\
Encapsulation & VXLAN/GENEVE streams require tunnel termination before sensor processing. &
  \texttt{isEncapsulationHandled} \\
Normalization & Canonical field roles are schema-specific; OCSF file-activity UIDs are computed
  and distinct. & \texttt{sourceIpFieldCode}, \texttt{fileSystemActivityTypeUid},
  \texttt{isKnownFileActivityUid} \\
Rule 1 & Administrative share plus write/create/open action plus executable extension is detected. &
  \texttt{matchesRemotePayloadRule} \\
Rule 2 & Default share observation with at least 15 distinct targets in 10 seconds is detected. &
  \texttt{matchesShareEnumerationRule} \\
Rule 3 & Rename or ransomware extension with at least 50 events in 60 seconds is detected. &
  \texttt{matchesRansomwareRule} \\
Rule 4 & Read action with at least 500 events in 300 seconds is detected. &
  \texttt{matchesExfiltrationRule} \\
Baseline tuning & Scanner/backup suppression blocks dispatch; untrusted off-hours context
  escalates to critical. & \texttt{shouldDispatchAlert}, \texttt{alertLevelForScore} \\
Correlation and SOAR & Endpoint correlation is required for dispatch; critical alerts select host
  isolation. & \texttt{hasEndpointCorrelation}, \texttt{responseForLevel} \\
\bottomrule
\end{tabularx}
\caption{Traceability from planning requirements to LexLean declarations.}
\end{table}

\section{Verification Evidence}

The full pipeline was executed from the project root:

\begin{lstlisting}
lexlean fmt --all --check
lexlean lock --check
lexlean check --all
lexlean verify --all
lake build --allow-empty
prismpm check
prismpm build
\end{lstlisting}

Recorded results after this integration:

\begin{itemize}
  \item Source modules: 14 (nine pre-existing plus five new enterprise modules).
  \item Theorems: 52; definitions: 99; structures: 39; inductive types: 24.
  \item LexLean semantic ID:
        \texttt{82a79bd033b89dc60f05bd8f58518b2f5d77d6e10c7b839ae02296ab422611b2}.
  \item Lean attestation:
        \texttt{d633598a6e4acdde2a11a1ddc39e35eaeb3d8a2d6d61510c5508b266671b2460}.
  \item PrismPM release:
        \texttt{79af3626858c67b10dad330a8e0a4ece40b728d9cde4072af3c9fe34da4a92b3}.
  \item The prior certified Rust production export
        \texttt{a69be6edb75c35ed99b44712e5a8e986b45586515bb78ddaeda13c0b4d9b2a9b} remains valid
        because its root and source module were not changed by this ADR.
\end{itemize}

\section{Consequences}

\subsection{Positive}

\begin{itemize}
  \item The planning document's structural and analytical claims are now present as named,
        checkable formal objects rather than prose only.
  \item Threshold drift is visible in source: changing a documented count or window changes a
        theorem's meaning and requires re-verification.
  \item The existing Wireshark suite and the new enterprise suite are jointly checked by
        \texttt{Main}, giving a single project-wide attestation.
  \item PrismPM produces a fresh content-addressed release for the expanded project.
\end{itemize}

\subsection{Negative and Mitigations}

\begin{itemize}
  \item The model uses finite codes for IP fields, shares, actions, extensions, and schemas rather
        than raw strings or packet bytes. \textbf{Mitigation:} this is intentional bounded
        formalization; no parsing equivalence is claimed.
  \item Operational systems such as Kafka, Zeek, Suricata, Elasticsearch, and SOAR are not deployed
        or executed. \textbf{Mitigation:} downstream engineering must treat this ADR as a
        specification of decision structure, not as a deployment artifact.
  \item Statistical baselines are represented by discrete weights and classes. \textbf{Mitigation:}
        the formal model proves threshold behavior for representative contexts; empirical tuning
        remains outside scope.
\end{itemize}

\section{Status}

Accepted and implemented on 2026-11-23. Any future change to the modeled thresholds, schema codes,
or tier ordering requires a new ADR or an explicit amendment to this record, followed by a complete
LexLean verification and PrismPM rebuild.

\end{document}
